\documentclass[conference]{IEEEtran}

\usepackage[utf8]{inputenc}
\usepackage[T1]{fontenc}
\usepackage{amsmath,amssymb}
\usepackage{booktabs}
\usepackage{tabularx}
\usepackage{cite}
\usepackage[hidelinks]{hyperref}

\newcommand{\SO}{\mathrm{SO}(3)}
\newcommand{\SE}{\mathrm{SE}(3)}

\begin{document}

\title{M1: EquiGraspFlow Baseline\\[0.3em]
  {\normalsize Milestone report $\cdot$ ADLR WS 2026/27, project idea 4}}

\author{
  \IEEEauthorblockN{Georgii Fatakhov}
  \IEEEauthorblockA{Technical University of Munich}
  \and
  \IEEEauthorblockN{Awar Satar}
  \IEEEauthorblockA{Technical University of Munich}
}

\maketitle

\section{Goal and Setup}

M1 establishes the flow-matching baseline: evaluate the released EquiGraspFlow~\cite{lim2024equigraspflow}
checkpoints, then train our own model, whose EMD should be close to the published value and to the
released checkpoint. All runs use the unmodified reference code.

\textbf{Data.} ACRONYM~\cite{eppner2021acronym} subset released with EquiGraspFlow: Laptop, Mug, Bowl and
Pencil, with $105/35/35$, $54/18/22$, $21/7/10$ and $50/16/16$ objects for train/val/test and one model
for all categories. All 449 grasp files and meshes load. The full point cloud has 1024 points sampled on the mesh, the partial one 512
points visible from one viewpoint.

\textbf{Protocol.} Each test object is presented in 3 evenly spaced rotations (partial: $\times\,3$
viewpoints), giving 249 (partial: 747) test cases. Sampling uses the 4th-order Runge--Kutta--Munthe-Kaas
solver on $\SE$ with 20 steps, i.e.\ 80 evaluations of the guided vector field, each running the network
twice for classifier-free guidance.

\textbf{Metric.} For one test case the model generates as many grasps as the object has ground-truth
grasps (about 1.1--1.9k). The EMD is the mean pair distance $\sqrt{\theta^2 + \lVert \Delta p \rVert^2}$
under the optimal one-to-one assignment, with $\theta$ the relative rotation angle in rad and $\Delta p$
the translation difference in m. It is averaged over the test cases of each category.

\section{Results}

Table~\ref{tab:emd} compares the released checkpoints with the paper. All categories agree within
$0.003$ except Mug with full point clouds ($+0.014$); a different random stream changes a category mean
by about $0.005$. The full test set takes 248\,s (partial: 500\,s) on an RTX 5070 Ti.

\begin{table}[!t]
  \centering
  \caption{Test EMD (lower is better). Paper: SO(3)-aug rows of Tables 1a and 5a in~\cite{lim2024equigraspflow}.}
  \label{tab:emd}
  \begin{tabular}{@{}llcccc@{}}
    \toprule
    Point cloud & Model & Laptop & Mug & Bowl & Pencil \\
    \midrule
    full & Paper & 0.342 & 0.483 & 0.313 & 0.287 \\
    full & Released ckpt. & 0.340 & 0.497 & 0.314 & 0.286 \\
    full & Own training & \multicolumn{4}{c}{pending} \\
    \midrule
    partial & Paper & 0.303 & 0.409 & 0.268 & 0.265 \\
    partial & Released ckpt. & 0.301 & 0.412 & 0.267 & 0.263 \\
    \bottomrule
  \end{tabular}
\end{table}

\begin{table}[!t]
  \centering
  \caption{EMD reference floor: two random halves of the ground-truth grasps, and model samples against
  one half (full point cloud, released checkpoint, half the grasps per side).}
  \label{tab:floor}
  \begin{tabular}{@{}lcccc@{}}
    \toprule
    & Laptop & Mug & Bowl & Pencil \\
    \midrule
    Ground truth vs.\ ground truth & 0.253 & 0.298 & 0.222 & 0.262 \\
    Model vs.\ ground truth & 0.368 & 0.545 & 0.337 & 0.326 \\
    \bottomrule
  \end{tabular}
\end{table}

\textbf{Interpreting the EMD.} The EMD has no zero point in practice: two finite samples of the same grasp
distribution are already far apart, as Table~\ref{tab:floor} shows. The distance is dominated by
rotation, and grasps that differ by the $180^\circ$ symmetry of the gripper count as $\pi$ apart. The
model is closest to the floor on Pencil, Bowl and Laptop; Mug shows the largest gap, consistent with it
having the highest EMD in the paper. EMD values are therefore compared with each other at equal sample
size rather than read as absolute quality; physical success is measured in MuJoCo in M2.

\section{Deviations and Notes}

\begin{itemize}
  \item The released configs use weight decay $10^{-6}$, the paper text states $10^{-5}$; we follow the
  configs. The partial config uses guidance $1.5$, the paper states $2$.
  \item On Windows the data loader must run without worker processes, since the dataset holds open3d
  meshes that cannot be pickled.
  \item Training with the stock config (40{,}000 epochs, about 2.3M iterations) takes about 0.21\,s per
  iteration locally, roughly 5.5 days. The released full checkpoint reached its best validation loss at
  iteration 700k.\end{itemize}

\section{Next Steps}

Train our own full-point-cloud model and evaluate it with the same protocol to complete
Table~\ref{tab:emd}. The evaluation also accepts other samplers and step counts, as needed for the NFE
sweeps of M2 and M3~\cite{kwon2026graspmeanflow}.

\bibliographystyle{IEEEtran}
\bibliography{adlr}

\end{document}
